\section{Formal Estimand Definitions}
\label{sec:estimand-formalization}

Let reviewer \(r\) receive prompt \(T_r\in\{0,1\}\), rate item \(i\), and emit
label \(Y_{ri}(T_r)\). Let \(Z_i\) denote a constructed reference and
\(s(Y,Z)\) a reference-relative score.

\paragraph{Realized-panel contrast.}
\[
\begin{aligned}
\Delta_{\mathrm{panel}}
= \frac{1}{I}\sum_i \bigl[&
s\!\left(M_{r\in\mathcal R_1}Y_{ri}(1),Z_i\right)\\[-2pt]
&-s\!\left(M_{r\in\mathcal R_0}Y_{ri}(0),Z_i\right)
\bigr],
\end{aligned}
\]
where \(M\) is a majority rule and \(\mathcal R_t\) is the realized reviewer
panel. This quantity is descriptive for those panels and that reference.

\paragraph{Reviewer-population contrast.}
\[
\Delta_{\mathrm{rater}}
=
\mathbb E_{r,i,Z}\!\left[
s(Y_{ri}(1),Z_i)-s(Y_{ri}(0),Z_i)
\right].
\]
This target requires treatment assignment and uncertainty at the reviewer
level, plus a defined reviewer population and reference distribution.

\paragraph{Consequence contrast.}
\[
\Delta_{\mathrm{loss}}
=
\mathbb E\!\left[
L\{a(Y(1)),\theta\}-L\{a(Y(0)),\theta\}
\right],
\]
where \(a(\cdot)\) is the accept, narrow, reject, or defer decision and
\(\theta\) is the deployment state.

\paragraph{Variance decomposition.}
Let \(\mathcal R\) be the realized reviewer composition and \(\mathcal Z\) the
realized reference process. Then
\[
\begin{aligned}
\mathrm{Var}(\widehat\Delta)
={}&
\mathbb E\!\left[
\mathrm{Var}(\widehat\Delta\mid\mathcal R,\mathcal Z)
\right]\\
&+
\mathrm{Var}\!\left[
\mathbb E(\widehat\Delta\mid\mathcal R,\mathcal Z)
\right].
\end{aligned}
\]
Source resampling with reviewers and references fixed estimates the first
term. Reviewer assignment and reference construction contribute to the
second.

\paragraph{Assignment asymptotics for panel majorities.}
For three conditionally independent votes with correctness probabilities
\(p_1,p_2,p_3\), majority correctness has expectation
\[
m(p_1,p_2,p_3)
=p_1p_2+p_1p_3+p_2p_3-2p_1p_2p_3.
\]
Let \(Z\) be a six-versus-six reviewer assignment and
\(\mu(Z)=\mathbb E(\widehat\Delta\mid Z)\) the expected difference between
the two scheduled panel-majority scores. With \(I\) conditionally independent
items (or \(G\) independent source clusters of bounded size),
\[
\mathrm{Var}(\widehat\Delta)
=O(I^{-1})+\mathrm{Var}_{Z}\{\mu(Z)\}.
\]
The source-only standard error targets the first term. If reviewer profiles
are heterogeneous, the nonlinear function \(m\) generally makes
\(\mu(Z)\) nonconstant, so the second term does not vanish when more items are
added to the same twelve assignments. Conditional item precision can
therefore converge to zero while uncertainty over reviewer assignment remains.
This characterization explains the direction of the simulation; it does not
supply a population model for reviewer abilities.
